% ============================================================
% SECTION 5 — Imam al-Hassan al-'Askari (Chapter 7)
% ============================================================
\section{al-'Askari}

\begin{frame}{Imamate Under Total Surveillance}
\begin{itemize}[itemsep=5pt]
  \item Forced residence in Samarra's military quarter -- surveillance built into the environment itself
  \item Entire six-year Imamate lived either imprisoned or barred from normal social contact
  \item Surveillance intensifies sharply once the Mahdi's expected lineage becomes the specific concern
\end{itemize}
\bigquote{[Comparable to] Pharaoh's attempt to prevent the foretold birth among the Bani Isra'il.}{Mutahhari's comparison}
\end{frame}

\begin{frame}{Concealment That Survived the Searches}
\begin{itemize}[itemsep=5pt]
  \item Repeated, intensive searches for the promised child -- intensifying around the Imam's own death
  \item The birth successfully concealed from the general public -- known only to a small, trusted circle
  \item After his death: women of the household examined for pregnancy; one slave girl detained a full year on a mistaken suspicion
\end{itemize}
\delivernote{Invite a real, honest question here rather than rushing past it: given everything just described about the intensity of surveillance, how plausible is years-long successful concealment? What would it actually require logistically?}
\end{frame}

\begin{frame}{Jaddah and Hakimah Khatun: A Karbala Parallel}
\begin{itemize}[itemsep=5pt]
  \item Hudayth (``Jaddah'') -- standing rests on her \emph{own} learning, not merely on being al-'Askari's mother
  \item Becomes the community's learned refuge for religious questions after his death
  \item Hakimah Khatun draws the parallel explicitly: outward affairs to a trusted visible figure, while the true successor remains protected
\end{itemize}
\begin{center}
\begin{tikzpicture}[every node/.style={font=\scriptsize,align=center},
  box/.style={rectangle,rounded corners,draw=deepteal,fill=lightbg,text width=3.6cm,minimum height=1cm}]
\node[box] (a) {Karbala: Zaynab manages outward affairs};
\node[box, right=1.4cm of a] (b) {al-'Askari: Jaddah manages outward affairs};
\node[box, below=0.6cm of a] (c) {True successor: Zayn al-'Abidin, protected};
\node[box, below=0.6cm of b] (d) {True successor: hidden Imam al-Mahdi, protected};
\draw[-{Latex},thick,gold] (a) -- (c);
\draw[-{Latex},thick,gold] (b) -- (d);
\draw[dashed,thick,softgray] (a) -- (b) node[midway,above,font=\scriptsize\itshape]{same structure};
\end{tikzpicture}
\end{center}
\end{frame}

\qaframe{al-'Askari}{%
  \qa{Does the Pharaoh comparison do any real historical explanatory work, or is it purely theological?}
     {Both: theologically it reassures that state power cannot overturn what is divinely intended; historically it also explains \emph{why} the concealment strategy (trusted intermediaries, a small circle of knowledge) was genuinely necessary, not merely a passive trust in providence.}
  \qa{What does the level of scrutiny described here suggest about how carefully a reader should evaluate Chapter 8's later Mahdi-candidate claims?}
     {It raises the standard of evidence considerably -- if the genuine Mahdi's birth required this level of concealment to survive real state surveillance, Chapter 8's warnings about false claimants and mistaken allegiances carry extra weight.}
  \qa{Is Hakimah Khatun's Karbala/al-'Askari parallel doing the same interpretive work as the essence/expression principle from the Introduction?}
     {Yes -- both separate an outward, circumstance-driven arrangement (who visibly manages affairs) from the underlying, unchanged reality of who actually holds religious authority.}
}
